\section{Introduction}
\subsection{Page replacement algorithms}
\paragraph{First in First out}
What page replacament algorithm is being used? \\
We use FIFO (First in First out) page replacement algorithm for those laboratories as indicated
by the PageFault.java file line 18
\begin{figure}[H]
\caption{PageFault.java file}
\begin{lstlisting}[language=Java]
[...]
public class PageFault {

  /**
   * The page replacement algorithm for the memory management sumulator.
   * This method gets called whenever a page needs to be replaced.
   * <p>
   * The page replacement algorithm included with the simulator is 
   * FIFO (first-in first-out).  A while or for loop should be used 
   * to search through the current memory contents for a canidate 
   * replacement page.  In the case of FIFO the while loop is used 
   * to find the proper page while making sure that virtPageNum is 
   * not exceeded.
[...]
\end{lstlisting}
\end{figure}
All pages are stored in memory in a queue. Oldest page (First that came in) is
in front of this queue. \\ When we need to replace the page we remove the page that
is first in queue (so the one that came in as a first one, first in, first out)
\\
It is easy to explain and implement but in practical application it performs
poorly. It is still used but usually we use modified version of it.
\cite{Page Replacement Algorithms}
\paragraph{Optimal Page Replacement}
We replace pages which in the future will not be used for the longest time.
This is purely theoretical algorithm. It is perfect but not doable in practice
since operating systems can not know future requests. \\
It is used as a benchmark against which we compare other algorithms. 
\cite{Page Replacement Algorithms}
\paragraph{Least Recently Used}
We replace page which was not used for the longest time. \\
It is based on the idea that we will in future work on pages which we used
heavily in the past. In theory its performance can be close to optimal one but
in practice it is expensive to implement. \\
Most expensive way of implementing this algorithm is using linked lists. Most
recently used pages in front, least recently used in the back. Every time we
reference memory we have to move elements in list which takes a lot of time. \\
We can also use operating system counter \\
This algorithm is often used in different cheaper modified versions.
\cite{Page Replacement Algorithms}
\cite{pageWiki}
\subsection{Other}
\paragraph{Memory Management Unit} 
Hardware unit which translates virtual memory to physical one.
\cite{mmuWiki}
\paragraph{Page fault}
Exception raised when process wants to access page without preparations.
Preparation consists of adding the mapping to process's virtual address space
and/or loading page from a disk. MMU detects the fault but it is up to kernel
and/or loading page from a disk. 

